\documentclass[11pt]{article}

\usepackage[a4paper,margin=30mm]{geometry}
\usepackage[T1]{fontenc}
\usepackage[utf8]{inputenc}
\usepackage{lmodern}
\usepackage{microtype}
\usepackage{amsmath,amssymb,amsthm,mathtools}
\usepackage{booktabs}
\usepackage{enumitem}
\usepackage{array}
\usepackage{xurl}
\usepackage{hyperref}
\usepackage[nameinlink,noabbrev]{cleveref}

\hypersetup{
  pdftitle={A 2/3 Lower Bound for Progression-Free Harmonic Sums in Terms of van der Waerden Numbers},
  pdfauthor={Tim Held},
  pdfsubject={Erdos Problem 169; arithmetic progressions; van der Waerden numbers; Kempner sets},
  pdfkeywords={arithmetic progressions, van der Waerden numbers, harmonic sums, Kempner sets, additive combinatorics},
  colorlinks=true,
  linkcolor=black,
  citecolor=blue,
  urlcolor=blue
}
\urlstyle{same}

\newtheorem{theorem}{Theorem}[section]
\newtheorem{proposition}[theorem]{Proposition}
\newtheorem{lemma}[theorem]{Lemma}
\newtheorem{corollary}[theorem]{Corollary}
\theoremstyle{definition}
\newtheorem{definition}[theorem]{Definition}
\theoremstyle{remark}
\newtheorem{remark}[theorem]{Remark}

\newcommand{\N}{\mathbb{N}}
\newcommand{\HH}{\mathcal{H}}
\newcommand{\Kemp}{\mathcal{K}}
\newcommand{\diam}{\operatorname{diam}}

\title{A \texorpdfstring{$2/3$}{2/3} Lower Bound for Progression-Free Harmonic Sums\\
       in Terms of van der Waerden Numbers}
\author{Tim Held\\[2pt]
\small Independent researcher, Switzerland}
\date{Version 1.0, August 28, 2026}

\begin{document}
\maketitle

\begin{abstract}
For an integer $k\ge 3$, let $f(k)$ be the supremum of
$\sum_{n\in A}1/n$ over all sets $A$ of positive integers containing no
$k$-term arithmetic progression, and let $W(k)=W(2,k)$ be the two-color
van der Waerden number.  The classical coloring argument gives
$f(k)\ge \tfrac12 H_{W(k)-1}$ and hence
$\liminf_{k\to\infty} f(k)/\log W(k)\ge \tfrac12$.
We prove the stronger explicit estimate
\[
 f(k)\ge
 1+\frac{2W(k)-3}{3W(k)-5}
 \left(H_{W(k)-1}-\frac32\right),
\]
which implies
\[
 \liminf_{k\to\infty}\frac{f(k)}{\log W(k)}\ge \frac23.
\]
The proof starts with a two-coloring of $[W(k)-1]$ having no
monochromatic $k$-term progression.  Each finite color class is normalized
and used as the digit set of a shifted Kempner set in base $2W(k)-3$.
A modular progression argument proves that the resulting infinite set is
$k$-term-progression-free, while a leading-digit decomposition gives a
geometric amplification of its harmonic sum.  A weighted averaging step
then combines the two color classes.  We also give an $r$-color extension
and a coloring-sensitive refinement.  The result is a strict improvement
of the constant $1/2$, but it does not prove the stronger conjecture that
$f(k)/\log W(k)\to\infty$.
\end{abstract}

\noindent\textbf{Status and scope.}
This manuscript is a non-peer-reviewed preprint, and independent
verification is welcomed.  The principal result is a partial advance on
Erd\H{o}s Problem 169: it proves the weaker assertion that the classical
constant $1/2$ can be replaced by a larger absolute constant, while leaving
the stated limit-to-infinity problem open.  No claim is made here that the
full problem has been solved.

\medskip
\noindent\textbf{2020 Mathematics Subject Classification.}
Primary 11B25; Secondary 05D10, 11B05.

\noindent\textbf{Keywords.}
Arithmetic progressions, van der Waerden numbers, harmonic sums,
Kempner sets, digit-restricted sets, additive combinatorics.

\section{Introduction}

A $k$-term arithmetic progression, abbreviated $k$-AP, is a set of the form
\[
 \{a,a+d,\ldots,a+(k-1)d\},
 \qquad a,d\in\N,\ d>0.
\]
A set of positive integers is called \emph{$k$-AP-free} if it contains no
such progression.  Define
\begin{equation}\label{eq:def-f}
 f(k):=\sup\left\{
   \sum_{n\in A}\frac1n:
   A\subseteq\N\text{ is $k$-AP-free}
 \right\}.
\end{equation}
The finiteness of these quantities is tied to the Erd\H{o}s--Tur\'an
conjecture on arithmetic progressions \cite{erdos-turan-1936,gerver-1977}.
The present paper concerns only unconditional lower bounds, so no finiteness
assumption is used.

For integers $r\ge 2$ and $k\ge 2$, let $W(r,k)$ be the least positive
integer $N$ such that every coloring of $[N]:=\{1,\ldots,N\}$ with $r$
colors contains a monochromatic $k$-AP.  We write
\[
 W(k):=W(2,k).
\]
Van der Waerden's theorem guarantees that these numbers exist
\cite{vanderwaerden-1927}.  Erd\H{o}s and Graham asked for estimates on
$f(k)$ in terms of $W(k)$ and observed the elementary inequality
\begin{equation}\label{eq:classical-half-intro}
 f(k)\ge \frac12 H_{W(k)-1},
\end{equation}
where $H_N:=\sum_{n=1}^N1/n$ is the $N$th harmonic number.  In particular,
\[
 \liminf_{k\to\infty}\frac{f(k)}{\log W(k)}\ge\frac12.
\]
They explicitly asked whether the constant $1/2$ could be increased, and
more strongly whether
\begin{equation}\label{eq:erdos-question}
 \lim_{k\to\infty}\frac{f(k)}{\log W(k)}=\infty.
\end{equation}
See \cite{erdos-graham-1979,erdos-graham-1980,erdos-problems-169}.
Berlekamp obtained the linear lower bound
\[
 f(k)\ge \frac{\log 2}{2}\,k,
\]
and Gerver later proved the stronger asymptotic estimate
\begin{equation}\label{eq:gerver}
 f(k)\ge (1-o(1))k\log k
\end{equation}
\cite{berlekamp-1968,gerver-1977}.  The estimate in this paper addresses a
different scale: it specifically improves the coefficient attached to
$\log W(k)$.

Our main theorem is the following.

\begin{theorem}[Main theorem]\label{thm:main}
For every integer $k\ge 3$,
\begin{equation}\label{eq:main-exact}
 f(k)\ge
 1+\frac{2W(k)-3}{3W(k)-5}
 \left(H_{W(k)-1}-\frac32\right).
\end{equation}
Consequently,
\begin{equation}\label{eq:main-asymptotic}
 \liminf_{k\to\infty}\frac{f(k)}{\log W(k)}\ge\frac23.
\end{equation}
\end{theorem}

The original derivation of the argument gave the slightly weaker but still
sufficient estimate
\begin{equation}\label{eq:advertised-bound}
 f(k)\ge
 1+\frac{2W(k)-3}{3W(k)-5}
 \left(H_{W(k)-1}-2\right).
\end{equation}
Keeping track of the distinct minima of the two color classes replaces the
constant $2$ by $3/2$, producing \cref{eq:main-exact}.

The mechanism is simple.  Begin with a two-coloring of
$[W(k)-1]$ having no monochromatic $k$-AP.  Normalize a color class $C$ by
subtracting its minimum, and use the resulting finite set as the allowed
base-$b$ digits of a Kempner set.  If $b$ is greater than twice the diameter
of $C$, ordinary $k$-AP-freeness implies modular $k$-AP-freeness.  A
base-$b$ descent then shows that the infinite digit-restricted set is
$k$-AP-free.  Its successive digit layers contribute a geometric series to
the harmonic sum, with ratio $|C|/b$.  Applying this construction to both
color classes and averaging the two resulting bounds yields the constant
$2/3$.

Digit-restricted harmonic series originate with Kempner's classical example
\cite{kempner-1914}.  The relevant facts about progression-free digit sets
are closely related to the framework developed by Walker and Walker
\cite{walker-walker-2020} and by Walker \cite{walker-2025}.  Complete proofs
are included here so that the present argument is self-contained.

\section{Notation and elementary inequalities}

Throughout, $k\ge 3$ unless stated otherwise.  For a set
$A\subseteq\N$, write
\[
 \HH(A):=\sum_{a\in A}\frac1a
\]
with the value $+\infty$ permitted.  For a nonempty finite set $C$ of
integers, define
\[
 |C|=\text{its cardinality},\qquad
 \diam(C):=\max C-\min C.
\]
Translations preserve arithmetic progressions: $C$ is $k$-AP-free if and
only if $C+t$ is $k$-AP-free for every integer $t$.

We will repeatedly use a weighted averaging inequality.

\begin{lemma}[Weighted maximum]\label{lem:weighted-max}
Let $u_1,\ldots,u_q\in\mathbb{R}$ and let $d_1,\ldots,d_q>0$.  Then
\begin{equation}\label{eq:weighted-max}
 \max_{1\le i\le q}\frac{u_i}{d_i}
 \ge
 \frac{\sum_{i=1}^q u_i}{\sum_{i=1}^q d_i}.
\end{equation}
\end{lemma}

\begin{proof}
The quotient on the right is the weighted average
\[
 \frac{\sum_i d_i(u_i/d_i)}{\sum_i d_i}
\]
of the numbers $u_i/d_i$, with positive weights $d_i$.  A weighted average
cannot exceed the largest entry.  No sign condition on the $u_i$ is needed.
\end{proof}

For later limit arguments, note also that
\begin{equation}\label{eq:W-elementary-growth}
 W(r,k)\ge k+1
 \qquad (r\ge2,\ k\ge3).
\end{equation}
Indeed, color the interval $[k]$ with at least two colors.  Its only
$k$-term arithmetic progression is the whole interval $[k]$, which is not
monochromatic.  Thus $W(r,k)>k$.

We also record the standard harmonic-number asymptotic
\begin{equation}\label{eq:harmonic-asymptotic}
 H_N=\log N+\gamma+O(N^{-1}),
\end{equation}
where $\gamma$ is the Euler--Mascheroni constant.

\section{The classical \texorpdfstring{$1/2$}{1/2} bound}

The starting point is the short coloring argument of Erd\H{o}s and Graham.

\begin{proposition}\label{prop:classical-half}
For every $k\ge 3$,
\[
 f(k)\ge \frac12 H_{W(k)-1}.
\]
Consequently,
\[
 \liminf_{k\to\infty}\frac{f(k)}{\log W(k)}\ge\frac12.
\]
\end{proposition}

\begin{proof}
Put $N=W(k)-1$.  By the minimality of $W(k)$, there is a two-coloring
\[
 [N]=C_0\sqcup C_1
\]
with no monochromatic $k$-AP.  Thus both $C_0$ and $C_1$ are $k$-AP-free.
Moreover,
\[
 \HH(C_0)+\HH(C_1)=H_N.
\]
At least one of the two sums is at least $H_N/2$, and that color class is
admissible in the definition of $f(k)$.  This proves the finite inequality.
The asymptotic statement follows from \cref{eq:harmonic-asymptotic} and
$W(k)\to\infty$.
\end{proof}

The rest of the paper amplifies the harmonic sum of each finite color class
before the final averaging step.

\section{Modular arithmetic progressions and Kempner sets}

\subsection{Progressions modulo a base}

\begin{definition}\label{def:modular-free}
Let $b\ge 2$ and $S\subseteq\{0,1,\ldots,b-1\}$.  We say that $S$ contains
a \emph{$k$-term arithmetic progression modulo $b$} if there are integers
$x$ and $\Delta$ with $b\nmid\Delta$ such that
\[
 x+j\Delta\pmod b\in S
 \qquad (j=0,1,\ldots,k-1).
\]
If no such pair exists, $S$ is called \emph{$k$-AP-free modulo $b$}.
\end{definition}

The difference $\Delta$ need not be coprime to $b$, and the residues need
not be distinct.  This is the notion needed for the digit descent below.

\begin{proposition}[Large-base modular freeness]\label{prop:large-base}
Let $S\subseteq[0,M]$ be an ordinary $k$-AP-free set.  If $b>2M$, then
$S$ is $k$-AP-free modulo $b$.
\end{proposition}

\begin{proof}
Assume for contradiction that there are $x$ and $\Delta$ as in
\cref{def:modular-free}.  Replacing $\Delta$ by its least positive residue,
we may suppose
\[
 1\le \delta\le b-1,
 \qquad \delta\equiv\Delta\pmod b.
\]
Let
\[
 r_j\in S\subseteq[0,M]
\]
be the least nonnegative residue of $x+j\delta$ modulo $b$.  Since
$0<\delta<b$, each successive difference is either
\[
 r_{j+1}-r_j=\delta
 \quad\text{or}\quad
 r_{j+1}-r_j=\delta-b.
\]
If a step of the first type occurs, then $\delta\le M$, because both
residues lie in $[0,M]$.  If a step of the second type occurs, then
$b-\delta\le M$.

If both step types occur, then
\[
 b=\delta+(b-\delta)\le 2M,
\]
contrary to $b>2M$.  Hence every step has the same type.  If all steps are
$+\delta$, then $r_0,r_1,\ldots,r_{k-1}$ is an ordinary increasing $k$-AP
in $S$.  If all steps are $\delta-b$, then the same residues, read in reverse
order, form an ordinary increasing $k$-AP in $S$.  Both alternatives
contradict the assumed $k$-AP-freeness of $S$.
\end{proof}

\subsection{Digit-restricted sets}

\begin{definition}[Kempner set]\label{def:kempner}
Let $b\ge 2$ and let $S\subseteq\{0,1,\ldots,b-1\}$ with $0\in S$.  The
associated Kempner set is
\[
 \Kemp(S,b):=
 \left\{
   \sum_{j=0}^{L} a_jb^j:
   L\ge 0,\ a_j\in S
 \right\}.
\]
Equivalently, $\Kemp(S,b)$ is the set of nonnegative integers all of whose
canonical base-$b$ digits belong to $S$.
\end{definition}

\begin{theorem}[Digit descent]\label{thm:digit-descent}
If $S$ is $k$-AP-free modulo $b$ and $0\in S$, then $\Kemp(S,b)$ is
ordinary $k$-AP-free.  Consequently, $\Kemp(S,b)+1$ is a $k$-AP-free set of
positive integers.
\end{theorem}

\begin{proof}
Suppose, toward a contradiction, that
\[
 x,x+d,\ldots,x+(k-1)d
\]
is contained in $\Kemp(S,b)$, where $d>0$.  Let $t\ge 0$ be the largest
integer for which $b^t\mid d$.  We remove $t$ common terminal digits.

If $b\mid d$, all progression terms have the same units digit, say
$a_0\in S$.  After subtracting $a_0$ from every term and dividing by $b$,
we obtain another $k$-AP contained in $\Kemp(S,b)$, now with common
difference $d/b$.  Repeating this operation exactly $t$ times produces a
$k$-AP in $\Kemp(S,b)$ whose common difference $d'=d/b^t$ is not divisible
by $b$.

Reducing the final progression modulo $b$ gives
\[
 x'+jd'\pmod b\in S
 \qquad (j=0,\ldots,k-1),
\]
with $b\nmid d'$.  This is a $k$-term arithmetic progression modulo $b$ in
$S$, contradicting the hypothesis.  Therefore $\Kemp(S,b)$ is $k$-AP-free.
Translation by $1$ preserves arithmetic progressions, so
$\Kemp(S,b)+1$ is also $k$-AP-free.
\end{proof}

Combining \cref{prop:large-base,thm:digit-descent} gives the form used later.

\begin{corollary}[Finite seed to infinite set]\label{cor:finite-seed}
Let $C$ be a nonempty finite $k$-AP-free set of integers.  Put
\[
 m=\min C,
 \qquad
 S=C-m,
 \qquad
 M=\diam(C).
\]
For every integer $b>2M$, the shifted Kempner set
\[
 \Kemp(S,b)+1
\]
is a $k$-AP-free set of positive integers.
\end{corollary}

\begin{proof}
The translated set $S$ is ordinary $k$-AP-free, lies in $[0,M]$, and
contains $0$.  Apply \cref{prop:large-base,thm:digit-descent}.
\end{proof}

\section{Harmonic-sum amplification}

The next estimate is the quantitative component of the argument.

\begin{proposition}[Leading-digit lower bound]\label{prop:leading-digit}
Let $b\ge 2$, let $S\subsetneq\{0,1,\ldots,b-1\}$ contain $0$, and write
$s=|S|$.  Then
\begin{equation}\label{eq:leading-digit}
 \HH(\Kemp(S,b)+1)
 \ge
 1+\frac{\displaystyle
     \sum_{a\in S\setminus\{0\}}\frac1{a+1}}
          {\displaystyle 1-\frac{s}{b}}.
\end{equation}
In particular, the harmonic series of $\Kemp(S,b)+1$ converges whenever
$s<b$.
\end{proposition}

\begin{proof}
The element $0\in\Kemp(S,b)$ contributes $1$ after the shift by $1$.
Now fix an integer $\ell\ge 1$ and consider the positive elements of
$\Kemp(S,b)$ whose canonical base-$b$ expansion has exactly $\ell$ digits.
Each such integer has a unique representation
\[
 n=ab^{\ell-1}+q,
\]
where
\[
 a\in S\setminus\{0\}
\]
is the leading digit and $q$ is an integer represented by $\ell-1$ digits
from $S$, allowing leading zeroes in this lower block.  There are exactly
$s^{\ell-1}$ choices for $q$.

Since $0\le q\le b^{\ell-1}-1$, we have
\[
 n+1=ab^{\ell-1}+q+1
 \le (a+1)b^{\ell-1}.
\]
Therefore
\[
 \frac1{n+1}
 \ge
 \frac{1}{(a+1)b^{\ell-1}}.
\]
Summing first over all $q$, then over all nonzero leading digits $a$, the
contribution of the $\ell$-digit layer is at least
\[
 \left(\frac{s}{b}\right)^{\ell-1}
 \sum_{a\in S\setminus\{0\}}\frac1{a+1}.
\]
Summing these lower bounds over $\ell\ge 1$ and adding the contribution of
$0$ gives
\[
 \HH(\Kemp(S,b)+1)
 \ge
 1+
 \left(\sum_{a\in S\setminus\{0\}}\frac1{a+1}\right)
 \sum_{j=0}^{\infty}\left(\frac{s}{b}\right)^j.
\]
Because $S$ is a proper subset of the $b$ possible digits, $s<b$, so the
geometric series equals $(1-s/b)^{-1}$.  This proves
\cref{eq:leading-digit}.  For completeness, convergence follows from a separate upper estimate.
Every positive element in the $\ell$-digit layer is at least
$b^{\ell-1}$, while that layer contains at most $s^\ell$ elements.
Consequently its harmonic contribution after shifting is at most
$b(s/b)^\ell$.  The sum of these upper bounds is geometric because
$s<b$.
\end{proof}

For a finite seed $C$, it is convenient to introduce its normalized
harmonic weight
\begin{equation}\label{eq:phi-def}
 \Phi(C):=
 \sum_{c\in C}\frac1{c-\min C+1}.
\end{equation}

\begin{theorem}[Finite-seed amplification]\label{thm:seed-amplification}
Let $C$ be a nonempty finite $k$-AP-free set of positive integers, and let
$b$ be an integer satisfying
\[
 b>2\diam(C).
\]
Then
\begin{equation}\label{eq:seed-amplification}
 f(k)\ge
 1+\frac{\Phi(C)-1}{1-|C|/b}.
\end{equation}
\end{theorem}

\begin{proof}
Let $m=\min C$ and $S=C-m$.  By \cref{cor:finite-seed}, the set
$\Kemp(S,b)+1$ is $k$-AP-free.  Moreover,
\[
 \sum_{a\in S\setminus\{0\}}\frac1{a+1}
 =
 \sum_{c\in C\setminus\{m\}}\frac1{c-m+1}
 =\Phi(C)-1.
\]
Applying \cref{prop:leading-digit} and then the definition of $f(k)$ gives
\cref{eq:seed-amplification}.
\end{proof}

\begin{remark}\label{rem:shift-benefit}
For every nonempty finite $C\subseteq\N$ with $m=\min C$,
\begin{equation}\label{eq:phi-vs-h}
 \Phi(C)
 \ge
 \HH(C)+1-\frac1m.
\end{equation}
Indeed, the term $c=m$ changes from $1/m$ to $1$, and every other
denominator decreases from $c$ to $c-m+1\le c$.  The correction
$1-1/m$ is the source of the improvement from $H_N-2$ to $H_N-3/2$ in the
two-color theorem.
\end{remark}

\section{Amplifying a partition of an interval}

We first state a coloring-sensitive form that retains the individual class
sizes, spans, and normalized harmonic weights.

\begin{proposition}[Coloring-sensitive partition bound]\label{prop:partition-sensitive}
Suppose
\[
 [N]=C_1\sqcup\cdots\sqcup C_q
\]
is a partition into nonempty $k$-AP-free sets.  For each $i$, choose an
integer $b_i>2\diam(C_i)$ and put
\[
 s_i=|C_i|,
 \qquad
 d_i=1-\frac{s_i}{b_i}.
\]
Then
\begin{equation}\label{eq:partition-sensitive}
 f(k)\ge
 1+\frac{\displaystyle\sum_{i=1}^q\bigl(\Phi(C_i)-1\bigr)}
          {\displaystyle\sum_{i=1}^q\left(1-\frac{s_i}{b_i}\right)}.
\end{equation}
\end{proposition}

\begin{proof}
By \cref{thm:seed-amplification}, for every $i$,
\[
 f(k)\ge 1+\frac{\Phi(C_i)-1}{d_i}.
\]
The denominators $d_i$ are positive because
$s_i\le\diam(C_i)+1<b_i$.  Taking the maximum over $i$ and applying
\cref{lem:weighted-max} with $u_i=\Phi(C_i)-1$ proves
\cref{eq:partition-sensitive}.
\end{proof}

For a universal estimate depending only on $N$ and the number of classes,
we use a common safe base.

\begin{theorem}[Universal partition amplification]\label{thm:partition-universal}
Suppose
\[
 [N]=C_1\sqcup\cdots\sqcup C_q
\]
is a partition into $q\ge 1$ nonempty $k$-AP-free sets.  Then
\begin{equation}\label{eq:partition-universal}
 f(k)\ge
 1+\frac{H_N-H_q}{\displaystyle q-\frac{N}{2N-1}}
 =
 1+\frac{2N-1}{(2q-1)N-q}\,(H_N-H_q).
\end{equation}
\end{theorem}

\begin{proof}
Choose the common base
\[
 b=2N-1.
\]
For every class $C_i$,
\[
 \diam(C_i)\le N-1,
\]
so $b>2\diam(C_i)$.  Apply \cref{prop:partition-sensitive} with
$b_i=b$ for all $i$.  Since $\sum_i s_i=N$, its denominator is
\begin{equation}\label{eq:partition-denominator}
 \sum_{i=1}^q\left(1-\frac{s_i}{b}\right)
 =q-\frac{N}{b}
 =q-\frac{N}{2N-1}.
\end{equation}

It remains to lower-bound the numerator.  Let
\[
 m_i=\min C_i.
\]
By \cref{eq:phi-vs-h},
\[
 \Phi(C_i)-1
 \ge
 \HH(C_i)-\frac1{m_i}.
\]
Summing over the partition gives
\begin{equation}\label{eq:numerator-minima}
 \sum_{i=1}^q\bigl(\Phi(C_i)-1\bigr)
 \ge
 H_N-\sum_{i=1}^q\frac1{m_i}.
\end{equation}
The minima $m_1,\ldots,m_q$ are distinct positive integers.  The sum of
their reciprocals is therefore at most the sum of the reciprocals of the
$q$ smallest positive integers:
\begin{equation}\label{eq:minima-hq}
 \sum_{i=1}^q\frac1{m_i}\le H_q.
\end{equation}
Combining \cref{eq:partition-denominator,eq:numerator-minima,eq:minima-hq}
with \cref{eq:partition-sensitive} proves the first expression in
\cref{eq:partition-universal}.  Multiplying numerator and denominator by
$2N-1$ gives the second expression.
\end{proof}

\begin{remark}[The initially advertised estimate]\label{rem:advertised}
If one uses only the cruder inequality $\Phi(C_i)\ge\HH(C_i)$, then the
numerator in \cref{eq:partition-universal} is replaced by $H_N-q$.  For
$q=2$, this gives exactly
\[
 f(k)\ge 1+\frac{2N-1}{3N-2}(H_N-2),
\]
which is the bound in the original proof sketch.  The minima argument costs
nothing and improves $H_N-2$ to $H_N-H_2=H_N-3/2$.
\end{remark}

\section{Proof of the main theorem}

\begin{proof}[Proof of \cref{thm:main}]
Set
\[
 N=W(k)-1.
\]
By the definition of $W(k)$, there exists a two-coloring
\[
 [N]=C_0\sqcup C_1
\]
with neither color class containing a $k$-AP.  The two classes are nonempty:
if one were empty, the other would equal $[N]$, which contains $k$
consecutive integers and therefore a $k$-AP, since $N\ge k$.

Apply \cref{thm:partition-universal} with $q=2$.  Since $H_2=3/2$,
\[
 f(k)
 \ge
 1+\frac{2N-1}{3N-2}\left(H_N-\frac32\right).
\]
Substituting $N=W(k)-1$ yields
\[
 2N-1=2W(k)-3,
 \qquad
 3N-2=3W(k)-5,
\]
and therefore
\[
 f(k)\ge
 1+\frac{2W(k)-3}{3W(k)-5}
 \left(H_{W(k)-1}-\frac32\right).
\]
This is \cref{eq:main-exact}.

For the asymptotic consequence, note that $W(k)\ge k$, hence
$W(k)\to\infty$.  With $N=W(k)-1$, we have
\[
 \frac{2N-1}{3N-2}\longrightarrow\frac23,
 \qquad
 \frac{H_N-3/2}{\log(N+1)}\longrightarrow 1
\]
by \cref{eq:harmonic-asymptotic}.  Dividing the finite inequality by
$\log W(k)=\log(N+1)$ and taking the lower limit gives
\[
 \liminf_{k\to\infty}\frac{f(k)}{\log W(k)}\ge\frac23.
\]
\end{proof}

The new finite bound is not merely asymptotically stronger than the
classical coloring bound.  It is strictly stronger for every relevant
$N$.

\begin{corollary}[Exact gain over the classical bound]\label{cor:exact-gain}
Let $N=W(k)-1$.  Then the right-hand side of \cref{eq:main-exact} exceeds
$H_N/2$ by
\begin{equation}\label{eq:exact-gain}
 \frac{NH_N-1}{2(3N-2)}.
\end{equation}
In particular, the improvement is strict for every $k\ge 3$.
\end{corollary}

\begin{proof}
Write $a_N=(2N-1)/(3N-2)$.  Direct algebra gives
\begin{align*}
 \left[1+a_N\left(H_N-\frac32\right)\right]-\frac{H_N}{2}
 &=1-\frac32a_N+\left(a_N-\frac12\right)H_N\\
 &=\frac{NH_N-1}{2(3N-2)}.
\end{align*}
For $k\ge3$, one has $N\ge 2$, so $NH_N>1$.
\end{proof}

Since $H_N-3/2\ge H_N-2$, \cref{eq:main-exact} also immediately implies
the originally claimed bound \cref{eq:advertised-bound}.

\section{An \texorpdfstring{$r$}{r}-color extension}

The same argument applies to any number of colors.

\begin{theorem}[Van der Waerden extension]\label{thm:r-color}
Let $r\ge2$ and $k\ge3$, and put
\[
 N=W(r,k)-1.
\]
Then
\begin{equation}\label{eq:r-color}
 f(k)\ge
 1+\frac{2N-1}{(2r-1)N-r}\,(H_N-H_r).
\end{equation}
Consequently, for every fixed $r\ge2$,
\begin{equation}\label{eq:r-color-asymptotic}
 \liminf_{k\to\infty}
 \frac{f(k)}{\log W(r,k)}
 \ge
 \frac{2}{2r-1}.
\end{equation}
\end{theorem}

\begin{proof}
There is an $r$-coloring of $[N]$ without a monochromatic $k$-AP.  Discard
any empty color classes and let $q\le r$ be the number of nonempty classes.
By \cref{thm:partition-universal},
\[
 f(k)\ge
 1+\frac{H_N-H_q}{q-N/(2N-1)}.
\]
Because $q\le r$, we have $H_N-H_q\ge H_N-H_r$, while
$q-N/(2N-1)\le r-N/(2N-1)$.  Whenever $H_N-H_r\ge0$, replacing the
numerator by the smaller quantity and the denominator by the larger one
preserves a valid lower bound.  If $H_N-H_r<0$, the displayed target bound
is weaker than the trivial inequality $f(k)\ge1$ and is automatic.  Thus
\[
 f(k)\ge
 1+\frac{H_N-H_r}{r-N/(2N-1)},
\]
which is equivalent to \cref{eq:r-color}.

For fixed $r$, $W(r,k)\to\infty$ as $k\to\infty$.  Hence
$H_N-H_r\sim\log N$ and
\[
 \frac{2N-1}{(2r-1)N-r}\longrightarrow\frac{2}{2r-1}.
\]
Dividing by $\log W(r,k)=\log(N+1)$ proves
\cref{eq:r-color-asymptotic}.
\end{proof}

\section{Numerical illustrations}

The exact two-color van der Waerden numbers currently known are
\[
 W(3)=9,\qquad W(4)=35,\qquad W(5)=178,\qquad W(6)=1132.
\]
These values are recorded in OEIS A005346 \cite{oeis-a005346}.  The following
table compares three bounds obtained from them.  The second column of bounds is the classical coloring
estimate $H_N/2$.  The third is the initially advertised estimate with
$H_N-2$.  The fourth is the strengthened estimate of
\cref{eq:main-exact}, with $N=W(k)-1$.

\begin{table}[ht]
\centering
\caption{Illustrative finite lower bounds derived from exact values of $W(k)$.}
\label{tab:numerics}
\begin{tabular}{@{}c r r r r@{}}
\toprule
$k$ & $W(k)$ & $\tfrac12 H_N$ &
$1+\frac{2N-1}{3N-2}(H_N-2)$ &
$1+\frac{2N-1}{3N-2}(H_N-\tfrac32)$ \\
\midrule
3 & 9    & 1.358929 & 1.489448 & 1.830357 \\
4 & 35   & 2.059105 & 2.419201 & 2.754201 \\
5 & 178  & 2.878094 & 3.506492 & 3.840140 \\
6 & 1132 & 3.804258 & 4.739561 & 5.072944 \\
\bottomrule
\end{tabular}
\end{table}

These values are included only to illustrate the finite formula.  For small
$k$, specialized constructions, including those of Gerver and later
Kempner-set searches, give stronger numerical lower bounds.  The point of
\cref{thm:main} is the uniform coefficient of $\log W(k)$.

\subsection{A concrete three-term-progression-free example}

The case $k=3$ gives a direct check of every part of the construction.  Since
$W(3)=9$, the interval $[8]$ admits the progression-free two-coloring
\[
 C_0=\{1,2,5,6\},
 \qquad
 C_1=\{3,4,7,8\}.
\]
Both normalized color classes are equal to
\[
 S=\{0,1,4,5\}.
\]
The universal base is $b=2\cdot8-1=15$, and $15>2\max S=10$.
Consequently
\[
 A=\Kemp(\{0,1,4,5\},15)+1
   =\{1,2,5,6,16,17,20,21,61,62,65,66,\ldots\}
\]
is $3$-AP-free.  Retaining the exact normalized weight of either color class
in \cref{thm:seed-amplification} gives
\begin{align*}
 \HH(A)
 &\ge
 1+\frac{\frac12+\frac15+\frac16}{1-\frac4{15}}\\
 &=1+\frac{13/15}{11/15}
 =\frac{24}{11}
 \approx 2.181818.
\end{align*}
This exceeds the universal value $1.830357$ in \cref{tab:numerics} because
the universal theorem deliberately discards coloring-specific information in
order to obtain a formula depending only on $W(k)$.

\section{Discussion, limitations, and possible refinements}

\subsection{What has and has not been proved}

The argument is unconditional and constructs explicit infinite
$k$-AP-free sets from any extremal van der Waerden coloring.  It proves a
strict improvement of the coefficient $1/2$ in the lower bound involving
$\log W(k)$.  It does not establish \cref{eq:erdos-question}; the ratio is
shown only to have lower limit at least $2/3$.

No assumption that $f(k)$ is finite is made.  If $f(k)=+\infty$ for some
$k$, all stated lower bounds remain valid.  The constructed Kempner sets
have convergent harmonic sums because their allowed-digit ratio satisfies
$|S|/b<1$.

\subsection{Why the constant \texorpdfstring{$2/3$}{2/3} appears}

In a balanced two-coloring of an interval of length $N$, a typical color
class has about $N/2$ elements.  The universally safe base in the argument
is approximately $2N$, because the elementary modular-freeness lemma
requires a base greater than twice the class diameter.  Thus the geometric
ratio in the harmonic amplification is approximately
\[
 \frac{|C|}{b}\approx\frac14,
\]
so a color-class harmonic sum is multiplied by approximately
\[
 \frac{1}{1-1/4}=\frac43.
\]
The original coloring supplies approximately one half of the interval's
harmonic mass.  Multiplying $1/2$ by $4/3$ gives $2/3$.

This heuristic explains the constant produced by the universal argument,
but it is not an optimality theorem.  Particular colorings may have smaller
class diameters, more favorable normalized harmonic weights, or digit sets
that are already progression-free modulo bases substantially smaller than
twice their diameter.  \Cref{prop:partition-sensitive} retains exactly this
information.

\subsection{Possible routes beyond the universal bound}

Any improvement beyond $2/3$ by this circle of ideas would need additional
input not used in the universal estimate.  Potential directions include:
\begin{enumerate}[label=(\roman*),leftmargin=2.4em]
 \item proving that some color class in every extremal two-coloring is
 progression-free modulo a base smaller than $2\diam(C)+1$;
 \item exploiting carries rather than excluding them through the condition
 $b>2\diam(C)$;
 \item combining information from both color classes inside a single
 digit system instead of constructing two separate Kempner sets;
 \item using the class-specific bases and normalized weights in
 \cref{prop:partition-sensitive}, together with structural information on
 extremal van der Waerden colorings.
\end{enumerate}
These are suggestions for further investigation, not claims that any of the
approaches must succeed.

\subsection{Relation to Gerver's bound}

Combining \cref{eq:gerver} and \cref{thm:main} gives the unconditional
summary
\[
 f(k)\ge
 \max\left\{
   (1-o(1))k\log k,
   \left(\frac23-o(1)\right)\log W(k)
 \right\}.
\]
The two terms measure different aspects of the problem.  Gerver's estimate
is direct in $k$, whereas the present estimate is designed to answer the
specific van der Waerden-number question posed by Erd\H{o}s and Graham.

\appendix

\section{Proof audit and edge cases}

This appendix isolates the logical dependencies and common points at which
a digit argument can fail.

\subsection{Canonical expansions and leading zeroes}

Every nonnegative integer has a unique canonical base-$b$ expansion, except
that arbitrary leading zeroes may be appended.  In
\cref{prop:leading-digit}, the leading digit $a$ is required to be nonzero,
so each positive integer is counted exactly once.  The lower block $q$ is
allowed to have leading zeroes, which gives exactly $s^{\ell-1}$ choices.

\subsection{Why modular freeness is necessary}

Ordinary $k$-AP-freeness of the digit set alone does not generally prevent
an arithmetic progression whose residues wrap around modulo $b$.  The
condition $b>2M$ in \cref{prop:large-base} prevents a mixture of wrapped and
unwrapped steps.  This is precisely where the factor of approximately two
in the chosen base enters.

\subsection{Termination of the descent}

The descent in \cref{thm:digit-descent} terminates because a positive
integer $d$ is divisible by only finitely many powers of $b$.  Primality of
$b$ is irrelevant.  The relevant exponent is simply the largest $t$ for
which $b^t\mid d$.

\subsection{Positivity of denominators}

In \cref{thm:seed-amplification}, $C$ lies in an interval of length
$\diam(C)+1$, so
\[
 |C|\le\diam(C)+1<b
\]
when $b>2\diam(C)$ and $\diam(C)\ge1$.  If $\diam(C)=0$, then $|C|=1$ and
one may take any $b\ge2$.  Thus $1-|C|/b>0$ in all cases.

\subsection{The weighted average with possibly negative numerators}

For a very small class, $\Phi(C_i)-1$ can be zero, and cruder substitutes
such as $\HH(C_i)-1$ can be negative.  \Cref{lem:weighted-max} remains valid
without a positivity assumption on the numerators.  Only the denominators
must be positive.

\subsection{No circular use of van der Waerden's theorem}

Van der Waerden's theorem is used only to know that $W(k)$ exists and that,
by minimality, an avoiding coloring of $[W(k)-1]$ exists.  The construction
of the infinite $k$-AP-free set is then verified directly by modular
arithmetic and digit descent.

\subsection{Algebraic check of the final coefficient}

With $N=W(k)-1$ and $b=2N-1$, the sum of the two geometric denominators is
\[
 2-\frac{N}{b}
 =2-\frac{N}{2N-1}
 =\frac{3N-2}{2N-1}.
\]
Taking the reciprocal gives
\[
 \frac{2N-1}{3N-2}
 =\frac{2W(k)-3}{3W(k)-5},
\]
which tends to $2/3$.

\section*{Declarations}

\noindent\textbf{Author responsibility and AI-assisted preparation.}
The manuscript was developed and typeset with substantial assistance from
GPT-5.6 Pro, used in the drafting session under the ``Sol Pro'' reasoning
configuration.  The system assisted with proof exploration, consistency
checks, literature organization, and LaTeX drafting.  Tim Held is the sole
human author, accepts responsibility for the contents, and is responsible
for independently verifying every claim before formal submission or
publication.  The AI system is not an author.

\medskip
\noindent\textbf{Funding.}
No external funding was received for this work.

\medskip
\noindent\textbf{Competing interests.}
The author declares no competing interests.

\medskip
\noindent\textbf{Data and code availability.}
No external dataset, computer-assisted proof, or custom verification code is
required for the argument.  The numerical values in \cref{tab:numerics} are
direct evaluations of the displayed formulas using the exact van der Waerden
numbers cited there.  The complete LaTeX source accompanies this preprint.

\begin{thebibliography}{99}
\small

\bibitem{berlekamp-1968}
E. R. Berlekamp,
\newblock A construction for partitions which avoid long arithmetic progressions,
\newblock \emph{Canadian Mathematical Bulletin} \textbf{11} (1968), no. 3,
409--414, doi:10.4153/CMB-1968-047-7.

\bibitem{erdos-problems-169}
Erd\H{o}s Problems,
\newblock Problem 169,
\newblock \url{https://www.erdosproblems.com/169}, accessed August 28, 2026.

\bibitem{erdos-graham-1979}
P. Erd\H{o}s and R. L. Graham,
\newblock Old and new problems and results in combinatorial number theory:
van der Waerden's theorem and related topics,
\newblock \emph{L'Enseignement Math\'ematique} (2) \textbf{25} (1979), no. 3--4,
325--344.

\bibitem{erdos-graham-1980}
P. Erd\H{o}s and R. L. Graham,
\newblock \emph{Old and New Problems and Results in Combinatorial Number Theory},
\newblock Monographies de L'Enseignement Math\'ematique, no. 28,
Universit\'e de Gen\`eve, 1980.

\bibitem{erdos-turan-1936}
P. Erd\H{o}s and P. Tur\'an,
\newblock On some sequences of integers,
\newblock \emph{Journal of the London Mathematical Society} \textbf{11} (1936),
261--264.

\bibitem{gerver-1977}
J. L. Gerver,
\newblock The sum of the reciprocals of a set of integers with no arithmetic
progression of $k$ terms,
\newblock \emph{Proceedings of the American Mathematical Society}
\textbf{62} (1977), no. 2, 211--214.

\bibitem{kempner-1914}
A. J. Kempner,
\newblock A curious convergent series,
\newblock \emph{The American Mathematical Monthly} \textbf{21} (1914), no. 2,
48--50.

\bibitem{oeis-a005346}
OEIS Foundation Inc.,
\newblock The On-Line Encyclopedia of Integer Sequences, A005346:
van der Waerden numbers $W(2,n)$,
\newblock \url{https://oeis.org/A005346}, accessed August 28, 2026.

\bibitem{vanderwaerden-1927}
B. L. van der Waerden,
\newblock Beweis einer Baudetschen Vermutung,
\newblock \emph{Nieuw Archief voor Wiskunde} \textbf{15} (1927), 212--216.

\bibitem{walker-2025}
A. Walker,
\newblock Integer sets of large harmonic sum which avoid long arithmetic
progressions,
\newblock arXiv:2203.06045v2 [math.NT], 2025,
\newblock \url{https://arxiv.org/abs/2203.06045}.

\bibitem{walker-walker-2020}
Aled Walker and Alexander Walker,
\newblock Arithmetic progressions with restricted digits,
\newblock \emph{The American Mathematical Monthly} \textbf{127} (2020), no. 2,
140--150, doi:10.1080/00029890.2020.1682888.

\end{thebibliography}

\end{document}
